\documentclass[11pt]{article}
\usepackage[margin=1in]{geometry}
\usepackage{amsmath,amssymb,amsthm}
\usepackage{graphicx}
\usepackage{microtype}
\usepackage[colorlinks=true,linkcolor=blue,citecolor=blue,urlcolor=blue]{hyperref}

\newtheorem{theorem}{Theorem}
\newtheorem{lemma}[theorem]{Lemma}
\theoremstyle{remark}
\newtheorem{remark}[theorem]{Remark}

\newcommand{\trN}{\operatorname{tr}_N}
\newcommand{\U}{\mathcal U}
\newcommand{\M}{\mathcal M}

\title{Two Explicit Witness Sectors of the\\
Random-Embedding Conjecture}
\author{Automated mathematical discovery run \texttt{fa\_banach\_001}\\
Lane 9, agent \texttt{agent\_lane\_09}, model GPT5.6}
\date{July 22, 2026}

\begin{document}
\maketitle

\begin{center}
\fbox{\parbox{0.9\textwidth}{\textbf{Status.}
Partial result likely valid; human review needed.  The full existential-
embedding conjecture remains open.}}
\end{center}

\section{Source conjecture}

Hayes, Jekel, and Kunnawalkam Elayavalli formulate two conjectures about
random embeddings of free group factors \cite{HJK}.  Their stronger
Conjecture 1.2, restated as Conjecture 5.1, is the following.  Let
$U_1^{(N)},\ldots,U_r^{(N)}$ be independent Haar random unitaries, and let
$u_1,\ldots,u_r$ be the canonical free Haar unitaries in $L(\mathbb F_r)$.
For every self-adjoint noncommutative $*$-polynomial $p$ in $r+s$ unitary
variables, they predict
\begin{equation}\label{eq:source-conjecture}
 \lim_{N\to\infty}\mathbb E\left[
 \sup_{V_1,\ldots,V_s\in\U(M_N(\mathbb C))}
 \trN p(U^{(N)},V)\right]
 =\sup_{v_1,\ldots,v_s\in\U(L(\mathbb F_r))}\tau p(u,v).
\end{equation}
They prove that this is equivalent to existentiality of the random embedding
of $L(\mathbb F_r)$ into a matrix ultraproduct.

\begin{figure}[ht]
\centering
\includegraphics[width=0.97\textwidth]{figures/open_problem_crop.png}
\caption{Conjecture 1.2 in the source paper, page 3.}
\end{figure}

We verify \eqref{eq:source-conjecture} for two infinite witness sectors.  The
first consists of all trace-affine witness polynomials.  The second consists
of separated terms in which one witness conjugates one self-adjoint base
polynomial.

\section{Proof intuition}

The witness optimization is difficult when several occurrences of the
auxiliary unitaries interact.  In the two sectors below, it has an exact
spectral solution before taking $N\to\infty$.

For a trace-affine term, polar decomposition gives
\[
             \sup_{V\in\U(M_N)}\operatorname{Re}\trN(CV)=\trN|C|.
\]
Thus the optimized formula depends only on the singular-value distribution
of a base polynomial.  For a single conjugation term, the rearrangement
inequality gives
\[
 \sup_{V\in\U(M_N)}\trN(AVBV^*)
 =\int_0^1\lambda_A(t)\lambda_B(t)\,dt,
\]
where $\lambda_A,\lambda_B$ are decreasing eigenvalue curves.  The same two
identities hold in the diffuse finite factor $L(\mathbb F_r)$.  Almost-sure
$*$-moment convergence of Haar matrices then gives convergence of trace
norms and of the decreasing eigenvalue curves.  This proves not only the
expected limit in the conjecture, but an almost-sure limit for every fixed
polynomial in either sector.

\section{Exact optimization lemmas}

Throughout, a finite factor carries its normalized trace $\tau$.  For a
self-adjoint $A$ in such a factor, let $\lambda_A:[0,1)\to\mathbb R$ denote
its decreasing eigenvalue rearrangement, or spectral scale.  For a matrix it
is the step function formed from the eigenvalues in decreasing order.

\begin{lemma}[Polar optimization]\label{lem:polar}
For every $C$ in a finite factor,
\begin{equation}\label{eq:polar}
                 \sup_{V\in\U(\M)}\operatorname{Re}\tau(CV)=\tau(|C|).
\end{equation}
\end{lemma}

\begin{proof}
The upper bound follows from noncommutative H\"older inequality.  Write
$C=W|C|$.  The polar partial isometry $W^*$ extends to a unitary $V$ because
its initial and final complements have the same trace.  For this extension,
$VW$ is the support projection of $|C|$, and cyclicity gives
$\tau(CV)=\tau(|C|VW)=\tau(|C|)$.
\end{proof}

\begin{lemma}[Rearrangement optimization]\label{lem:rearrangement}
If $A,B$ are self-adjoint elements of a diffuse finite factor, then
\begin{equation}\label{eq:rearrangement}
 \sup_{V\in\U(\M)}\tau(AVBV^*)
 =\int_0^1\lambda_A(t)\lambda_B(t)\,dt.
\end{equation}
\end{lemma}

\begin{proof}
The upper bound is the Hardy--Littlewood rearrangement inequality for finite
von Neumann algebras; it follows from the generalized singular-value
inequalities of Fack and Kosaki \cite{FK}.  The usual positive-operator form
implies the self-adjoint form after adding sufficiently large scalar
multiples of the identity to $A$ and $B$; all added terms are independent of
$V$.

For completeness, the reverse inequality can be seen directly.  Using
spectral scales, approximate $A$ and $B$ in $L_2$ by decreasing step
operators
\[
                 A_m=\sum_{k=1}^m a_kp_k,
       \qquad    B_m=\sum_{k=1}^m b_kq_k,
\]
where the $p_k$ and $q_k$ are orthogonal partitions of the identity with
$\tau(p_k)=\tau(q_k)=1/m$, and $a_k,b_k$ are decreasing.  Atoms of a spectral
measure may be split inside their spectral projections because the factor is
diffuse.  Projection families with the same trace vector are unitarily
conjugate, so there is $V_m$ with $V_mq_kV_m^*=p_k$ for every $k$.  Therefore
\[
             \tau(A_mV_mB_mV_m^*)=\frac1m\sum_{k=1}^ma_kb_k.
\]
The right side tends to the integral in \eqref{eq:rearrangement}.  The
$L_2$ approximation and Cauchy--Schwarz show that replacing $A_m,B_m$ by
$A,B$ changes the trace by $o(1)$, proving the lower bound.
\end{proof}

\section{The two witness sectors}

Let $q=q^*$ and all other capital letters below denote fixed
noncommutative $*$-polynomials in the base variables $X_1,\ldots,X_r$.

\begin{theorem}\label{thm:main}
For every fixed polynomial $p$ of either form below, the optimized random
matrix value converges almost surely to the $L(\mathbb F_r)$ value, and
\eqref{eq:source-conjecture} holds.

\medskip
\noindent\textup{(A) Trace-affine sector.}  For arbitrary $C_1,\ldots,C_s$,
\begin{equation}\label{eq:affine-form}
 p(X,Y)=q(X)+\sum_{j=1}^s
       \big(C_j(X)Y_j+Y_j^*C_j(X)^*\big).
\end{equation}
This includes, modulo cyclicity of the trace, every self-adjoint polynomial
of total witness degree at most one.

\medskip
\noindent\textup{(B) Separated one-conjugation sector.}  For self-adjoint
$A_j=A_j^*$ and $B_j=B_j^*$,
\begin{equation}\label{eq:conj-form}
 p(X,Y)=q(X)+\frac12\sum_{j=1}^s
 \big(A_j(X)Y_jB_j(X)Y_j^*
      +Y_jB_j(X)Y_j^*A_j(X)\big).
\end{equation}
\end{theorem}

\begin{proof}
We use Voiculescu's almost-sure asymptotic freeness theorem \cite{Voi}, which
gives, for every fixed $*$-polynomial $R$,
\begin{equation}\label{eq:moment-convergence}
 \trN R(U^{(N)})\longrightarrow\tau R(u)
 \qquad\text{almost surely}.
\end{equation}
All polynomial evaluations under consideration have a common uniform
operator-norm bound, since every variable is unitary.

For \textup{(A)}, Lemma~\ref{lem:polar}, applied independently to every
witness, gives the exact finite-$N$ formula
\begin{equation}\label{eq:affine-matrix}
 \sup_V\trN p(U^{(N)},V)
 =\trN q(U^{(N)})+2\sum_{j=1}^s\trN|C_j(U^{(N)})|.
\end{equation}
The identical factor formula is
\begin{equation}\label{eq:affine-factor}
 \sup_v\tau p(u,v)=\tau q(u)+2\sum_{j=1}^s\tau|C_j(u)|.
\end{equation}
By \eqref{eq:moment-convergence}, the constant terms converge.  To treat the
trace norms, apply \eqref{eq:moment-convergence} to powers of the positive
matrices
\[
                    C_j(U^{(N)})^*C_j(U^{(N)}).
\]
Uniform polynomial approximation of $t\mapsto\sqrt t$ on a common compact
interval yields
\[
                  \trN|C_j(U^{(N)})|\longrightarrow\tau|C_j(u)|.
\]
Equations \eqref{eq:affine-matrix} and \eqref{eq:affine-factor} prove the
almost-sure assertion for the affine sector.

For \textup{(B)}, cyclicity and Lemma~\ref{lem:rearrangement} give
\begin{align}
 \sup_V\trN p(U^{(N)},V)
 &=\trN q(U^{(N)})
   +\sum_{j=1}^s\int_0^1
     \lambda_{A_j(U^{(N)})}(t)\lambda_{B_j(U^{(N)})}(t)\,dt,
 \label{eq:conj-matrix}\\
 \sup_v\tau p(u,v)
 &=\tau q(u)
   +\sum_{j=1}^s\int_0^1
     \lambda_{A_j(u)}(t)\lambda_{B_j(u)}(t)\,dt.
 \label{eq:conj-factor}
\end{align}
Moment convergence and the common compact spectral bound imply weak
convergence of the empirical spectral measure of every self-adjoint base
polynomial to its spectral measure in $L(\mathbb F_r)$.  Generalized inverses
of uniformly compactly supported distribution functions therefore satisfy
\[
 \lambda_{A_j(U^{(N)})}(t)\longrightarrow\lambda_{A_j(u)}(t)
 \quad\text{for almost every }t,
\]
and likewise for $B_j$.  The eigenvalue curves are uniformly bounded, so
dominated convergence gives convergence of every integral in
\eqref{eq:conj-matrix} to its counterpart in \eqref{eq:conj-factor}.

Finally, in both sectors the optimized quantity is bounded uniformly in $N$
by a constant depending only on $p$.  Dominated convergence upgrades the
almost-sure limit to convergence of expectations, which is precisely
\eqref{eq:source-conjecture}.
\end{proof}

\section{Verification, scope, and novelty check}

The included numerical script constructs the polar and eigenvector-alignment
optimizers for $180$ random matrix pairs of orders $2$ through $10$.  Both
exact finite-dimensional formulas pass to tolerance $10^{-10}$.  This is a
smoke test only; the proof above is analytic.

The full conjecture remains open.  The theorem does not handle monomials with
mixed witnesses, higher witness degree, or simultaneous affine and
conjugation terms involving the same witness.  Those interactions are not
reduced to one-variable spectral data by the present argument.

Cheap run-index searches found no packet or attempt for arXiv:2606.02985.
Bounded primary searches on July 22, 2026 used the arXiv id, title, exact
Conjecture 1.2 label, ``existential random embedding,'' and trace-affine
unitary optimization.  They found the source and standard background on
trace rearrangements, but no later resolution or explicit record of these
two witness sectors.

\section{Human-review recommendation}

Review as a likely valid substantial partial result.  The affine sector is
elementary after the exact polar formula.  The highest-value checks are the
diffuse-factor lower bound in Lemma~\ref{lem:rearrangement}, and the use of
weak spectral convergence to pass to the eigenvalue-rearrangement pairing.
An expert should also assess whether these explicit sectors are already
regarded as folklore in the model theory/free probability community.

\begin{thebibliography}{9}
\bibitem{HJK}
B. Hayes, D. Jekel, and S. Kunnawalkam Elayavalli,
\emph{Questions on the structure of random embeddings of $L(\mathbb F_2)$},
arXiv:2606.02985v1, 2026.

\bibitem{FK}
T. Fack and H. Kosaki,
\emph{Generalized $s$-numbers of $\tau$-measurable operators},
Pacific J. Math. 123 (1986), no.~2, 269--300.

\bibitem{Voi}
D. Voiculescu,
\emph{Limit laws for random matrices and free products},
Invent. Math. 104 (1991), no.~1, 201--220.
\end{thebibliography}

\end{document}
